\section{Study 2: The Language of Punishment}

Study~1 found that punitiveness is most strongly tied to hostile dispositions---hatred, revenge, authoritarianism, racial resentment---and only weakly tied to concern about crime. Public talk about punishment looks different: officials, jurists, and ordinary citizens defending punitive policy almost always reach for prosocial reasons---deterrence, public safety, rehabilitation. So when our participants were asked to explain the sentences they had just imposed, would they reproduce the prosocial vocabulary of public discourse, or would the hostile substrate from Study~1 show through?

Study~2 takes up that question by analyzing the open-ended sentencing justifications participants wrote at the end of Study~1. Three accounts are on the table. \textit{Sincerity}: prosocial language is the genuine expression of prosocial motives, used most by people whose actual concern is preventing future harm. \textit{Individual facade}: prosocial language is a strategic mask, used most by hostile people to cover their actual motives. \textit{Cultural default}: prosocial language is the dominant available vocabulary for talking about punishment, so almost everyone uses it, regardless of psychology. The three accounts make different empirical predictions, which we test below.

\subsection{Method}

\subsubsection{Open-Ended Responses}

After choosing a sentence in Study~1, each participant was asked: ``In your own words, please explain why you recommended this sentence for Darryl. What was your reasoning?'' A follow-up prompt asked for elaboration. We concatenated the two responses for analysis ($M = 48.3$ words, $SD = 35.9$, $\mathit{Mdn} = 37$, range 5--231); all $N = 496$ participants wrote at least one response. The NLP pipeline ran in Python; the integration analyses linking text to the Study~1 scales ran in R. Code and data are at the project repository.

\subsubsection{Classification Methods}

We used five methods to characterize each response: a dictionary classifier, a zero-shot multi-label classifier, a forced-choice classifier, an embedding-based prototype-similarity scheme, and an unsupervised topic model. Every method ran on every response.

The \textit{dictionary classifier} scored each response against ten category dictionaries we built by hand: five prosocial themes (deterrence, incapacitation, rehabilitation, retribution, norm expression) and five dark themes (revenge, suffering, degradation, exclusion, victim focus). Counts were normalized by response length, and each response received the category with the highest normalized count. Seventy-five responses had no matches in any category and were left unclassified.

The \textit{zero-shot classifier} used DeBERTa-v3-large fine-tuned for natural language inference \citep{he2023deberta}. We supplied eight candidate labels: deterrence and prevention, incapacitation and public safety, rehabilitation and reform, proportional justice, societal condemnation, punishment and suffering, revenge and payback, and victim closure.\footnote{The DeBERTa and Claude classifiers used a label set in which deterrence and incapacitation were combined into one ``deterrence and prevention'' label, reflecting how those goals usually co-occur in everyday talk. The dictionary and embedding methods kept them as separate categories.} Each response received a probability for each label; the highest-probability label became its primary classification.

The \textit{forced-choice classifier} sent the same eight labels to a large language model (Claude Haiku~4.5; \citealp{anthropicclaude2024}) under a structured prompt asking for the single best-fitting label and a confidence score. This serves as a strong reference point because the model can read context that simpler classifiers miss.

The \textit{embedding method} encoded each response as a vector using three sentence-embedding models from different architectures and training corpora: BGE-large-en-v1.5 \citep{xiao2024bge}, voyage-3-large \citep{voyageai2024}, and Sentence-MPNet-base-v2 \citep{song2020mpnet, reimers2019sentence}. We then took the cosine similarity between each response embedding and seven prototype sentences---one per justification theme. Both the per-theme similarities and the prosocial-minus-dark gap were kept as continuous measures. We re-ran these analyses with five prototype sets to check that any finding did not turn on a single prototype wording (Supplementary Materials).

The \textit{topic model} was BERTopic \citep{grootendorst2022bertopic} run on the BGE embeddings: dimensionality reduction with UMAP \citep{mcinnes2018umap} followed by density-based clustering with HDBSCAN \citep{mcinnes2017hdbscan}. The model produced seven topic clusters and an outlier set of 131 responses (26.4\%)---a typical outlier rate for short opinion text. We also computed VADER sentiment compound scores \citep{hutto2014vader} and Empath category probabilities \citep{fast2016empath}, and kept word count as a control for response length.

\subsubsection{Predictions}

The three accounts make distinct predictions about how text features should correlate with the Study~1 scales. \textit{Sincerity} predicts a positive correlation between the prosocial--dark gap in language and the crime concerns cluster: people who really do worry about crime should produce more prosocial-sounding justifications. \textit{Individual facade} predicts a negative correlation between the prosocial--dark gap and hostile aggression: hostile people should produce more prosocial-sounding language as cover. \textit{Cultural default} predicts no reliable correlation between the prosocial--dark gap and either cluster. We tested these predictions using six different ways of measuring the prosocial--dark gap, so any conclusion does not rest on one methodological choice.

\subsubsection{Validation and Sensitivity Analyses}

We tested each classifier against a 40-sentence ground-truth set: ten clear examples for each of four core themes (deterrence, rehabilitation, retribution, revenge). Top-1 accuracy was 87.5\% for DeBERTa and 100.0\% for Claude. (A legacy BART-MNLI classifier scored 77.5\% and was retired from the pipeline.) We assessed cross-method convergence with pairwise raw agreement and Cohen's $\kappa$ on the prosocial/dark classification. To check prototype sensitivity, we re-ran the embedding analyses across five prototype sets and three embedding models---all 45 combinations are in the Supplementary Materials. To check vignette stability, we re-ran the central analyses within each Study~1 vignette and fit mixed-effects models with vignette as a random intercept.

\subsection{Results}

\subsubsection{The Prosocial Surface}

All five methods agreed on the surface: participants reached for prosocial themes much more often than dark themes. Across methods, between 70.8\% and 96.4\% of responses received a prosocial top category (Table~\ref{tab:nlp_method_summary}). No method assigned more than 30\% of responses to a dark category.

Which prosocial theme came out on top depended on how fine-grained the labels were. The two embedding methods that kept deterrence and incapacitation distinct were dominated by retribution-related categories (BGE: 41.9\%; voyage: 55.0\%); the dictionary classifier spread weight more evenly across retribution (21.2\%), rehabilitation (19.4\%), and incapacitation (17.7\%). The two classifiers that combined deterrence and incapacitation into one label were dominated by public safety (Claude: 38.1\%) and a punishment-and-suffering label (DeBERTa: 24.0\%, the highest dark-side share among the supervised methods).

Cross-method convergence was high on raw agreement ($.65$ to $.95$ across the ten method pairs) but modest on Cohen's $\kappa$ ($.02$ to $.31$). This is the $\kappa$ paradox under heavy class imbalance \citep{feinstein1990high}: when one class accounts for most of the data, even accurate classifiers can produce $\kappa$ near zero. The two embedding models with the most overlap, BGE and voyage, showed the highest $\kappa$ ($.31$). Full convergence statistics are in the Supplementary Materials.

\subsubsection{The Semantic Mismatch}

If language matched the Study~1 scales, the most punitive participants should sound the most aggressive---more talk of revenge, suffering, exclusion, and so on. They didn't. Pearson correlations between every text feature and every relevant Study~1 measure produced a 99-cell matrix (Figure~\ref{fig:facade_heatmap}; full table in Table~\ref{tab:nlp_facade}). Of the 33 cells that reached both significance ($p < .05$) and a meaningful effect size ($|r| \geq .10$), 32 were negative.

The participants scoring highest on punitiveness, hostile aggression, and authoritarianism did not sound most prosocial. Nor did they sound most openly dark. What they did, quietly, was suppress one specific prosocial theme: rehabilitation. BGE rehabilitation similarity correlated with punitiveness at $r = -.37$, more than twice the magnitude of any positive cell in the matrix. The Claude prosocial-minus-dark gap showed a parallel column of negative correlations with punitiveness, hostile aggression, and the suffering construct (all $r$s between $-.23$ and $-.28$).

The dark-theme similarities, by contrast, were flat: BGE suffering and revenge similarities both correlated with hostile aggression at $|r| < .03$. High-hostile participants did not endorse more revenge or more suffering. They just stopped talking about rehabilitation. The other prosocial themes (deterrence, incapacitation, retribution) and the dark themes themselves moved very little across the hostility distribution.

\subsubsection{Three Competing Explanations}

To distinguish the three accounts, we measured the prosocial--dark gap in six different ways and tested each against the predictions. Three measures used embeddings (BGE, voyage, mpnet aggregate prosocial-minus-dark similarities); three used classifiers (Claude, DeBERTa, and Dictionary aggregate prosocial-minus-dark counts). We assigned each measure a verdict: \textit{individual facade} if the gap correlated significantly negatively with hostile aggression and that magnitude exceeded the crime-concerns correlation; \textit{sincerity} if the gap correlated significantly positively with crime concerns; \textit{cultural default} otherwise. The full table is in Table~\ref{tab:nlp_multiverse}.

The verdicts split sharply by method type. Four of the six measures returned cultural-default verdicts. The three embedding-based gaps showed no reliable correlation with hostile aggression (BGE, voyage, and mpnet; all $|r| \leq .08$). The dictionary gap correlated with hostility ($r = -.10$, $p = .03$), but at the same magnitude as it did with crime concerns---failing the facade test, which requires an asymmetric relation. The two supervised-classifier gaps---Claude ML ($r = -.25$) and DeBERTa ML ($r = -.17$, both $p < .001$)---returned individual-facade verdicts. None returned a sincerity verdict.

The split between methods tells us something about what each one measures. The embeddings register overall vocabulary overlap, treating small wording differences as part of one continuous score; on that yardstick, high-hostile and low-hostile participants sound alike. The supervised classifiers instead assign each response to a single theme, and on that finer-grained measure, high-hostile participants do shift away from prosocial themes. But the shift is concentrated in the rehabilitation category alone, as the next analysis shows.

The theme-level correlations show this directly (Table~\ref{tab:nlp_theme_correlations}). Rehabilitation was consistently the most negative row, with correlations ranging from $r = -.15$ with hostile aggression to $r = -.37$ with punitiveness, and values of $-.20$ for RWA and $-.16$ for racial resentment in between. The other prosocial themes hovered near zero (all $|r| \leq .07$), and no dark theme correlated positively with any hostile measure. The rehabilitation-suppression pattern replicated across all three embedding models and within each Study~1 vignette.

\subsubsection{The Cultural Default Across the Political Spectrum}

One concern about the cultural-default reading is that liberals and conservatives might be using genuinely different language that happens to look similar once aggregated. The BERTopic analysis lets us check directly. Topic distribution differed reliably by both political group, $\chi^2(12) = 23.47$, $p = .024$, and hostile-aggression tertile, $\chi^2(12) = 21.38$, $p = .045$.

Both tests reach the conventional threshold, but the cell-level differences are small (Table~\ref{tab:nlp_political_moderation}). No topic appeared in fewer than 3\% or more than 31\% of responses in any subgroup, and the rank order of topic frequencies was identical across the three political groups. What we see is small differences against a shared rhetorical backdrop. Liberals, moderates, and conservatives reach for the same justification themes in roughly similar proportions, despite large mean-level differences in punitiveness and hostile aggression (conservatives $M_{\textrm{punit}} = .26$ vs.\ liberals $-.31$; conservatives $M_{\textrm{hostile}} = 4.12$ vs.\ liberals $3.33$).

\subsubsection{Text Features as Independent Predictors}

A final question: do the text features predict things the Study~1 scales miss? We tested this with two hierarchical regressions, one predicting the punitiveness aggregate and one predicting the (within-vignette $z$-scored) sentencing decision (Table~\ref{tab:nlp_regression}).\footnote{For the sentencing regression, we used the 8-item attitudinal punitiveness composite as the Step~1 predictor rather than the full punitiveness aggregate (which contains the sentencing decision as a component). Otherwise Step~1 $R^2$ would be artificially inflated and would mask the contribution of the text features.} Step~1 entered five Study~1 psychological predictors. Step~2 added four text features: BGE rehabilitation similarity, BGE suffering similarity, VADER sentiment, and word count.

The text features added meaningful prediction in both models. For the sentencing decision, $R^2$ rose from .121 to .255, a gain of $\Delta R^2 = .134$, $F(4, 486) = 21.84$, $p < .001$. For punitiveness, $R^2$ rose from .428 to .507, $\Delta R^2 = .079$, $F(4, 486) = 19.35$, $p < .001$.

The strongest text predictor in both models was BGE rehabilitation similarity. More rehabilitation-similar language predicted shorter sentences ($\beta = -6.91$) and lower punitiveness ($\beta = -4.39$, both $p < .001$), even after controlling for hostile aggression and the other psychological predictors. BGE suffering similarity ran in the opposite direction (sentencing $\beta = +3.15$; punitiveness $\beta = +1.84$; both $p < .001$). The rehabilitation effect held under a mixed-effects model with vignette as a random intercept ($\beta = -4.45$, $t = -8.42$, $p < .001$).

\subsection{Discussion}

Study~2 found something none of the three accounts predicts on its own. Almost every participant sounded prosocial: 70\% to 96\% of responses, depending on the method. But the participants who sounded most prosocial were not the ones lowest in hostile aggression, and the participants who sounded least prosocial were not the ones highest in dark themes. The most punitive participants didn't endorse more revenge or suffering; they quietly stopped talking about rehabilitation. Of the 33 cells in the facade matrix that reached both significance and a meaningful effect size, 32 were negative, and the strongest single one was the $-.37$ correlation between rehabilitation-themed language and punitiveness.

We call this pattern a cultural default with selective suppression. The prosocial vocabulary is widely shared. Liberals and conservatives, hostile and non-hostile, all reach for deterrence, incapacitation, and proportional justice in roughly equal measure. What differs across the hostility distribution is not whether prosocial language is invoked but \textit{which} prosocial theme is invoked. The most hostile and punitive participants drop rehabilitation and continue using the other prosocial themes at the same rate as everyone else.

This pattern is hard to square with a strong individual-facade view, which would predict that hostile people openly express dark themes once a method is sensitive enough to detect them. We found no positive correlations between dark-theme similarity and any psychological measure of hostility. The sincerity view also fails: it would predict that crime-concerned participants use more prosocial language, but that correlation was never positive across any of the six gap measures.

The hierarchical regressions point in a different direction: text features predicted what self-report scales did not ($\Delta R^2 = .134$ for sentencing; $\Delta R^2 = .079$ for punitiveness). The Study~1 scales and the text features capture different parts of the punishment mindset, with the rehabilitation-suppression effect showing up most clearly on the text side. Computational text analysis complements structured scales rather than replacing them: the language people use to justify punishment carries information the conventional instruments miss.

The General Discussion takes up these findings together with those of Study~1 to develop the argument that contemporary punishment discourse operates as a goal-mischaracterized practice: prosocial in its avowed purposes, but psychologically rooted in dispositions that those purposes do not name.
